% Separate research proof fragment. The open main manuscript is unchanged.
\subsection{Two disjoint fallback actions and feedback into the Chen costs}
\label{sec:third-prefix-feedback-transfer-20261009}

This consequence retains the original named distribution and ordinary
sieve inputs and their previously stated verification scopes. It uses
the own-route Chen lower count, not the external Li coefficient.
All constants refer to $H_N=C_NN/\log^2N$ with the existing singular
series and counting convention.

\paragraph{One common upper strategy.}
Let
\[
 E=\frac{2444205458396424207668787105}{361468820064774123237146624},
 \quad D_2=\frac{640339}{15210735000},
 \quad D_3=\frac{3097012678584315027795429}{2^{96}}.
\]
The second-prefix proof supplies the elementary lower bound $D_2$.
The third-prefix pair proof supplies $D_3$ by 4096 rational boxes.
Its complete independent replay and separate analytic review are
recorded in the companion review, not inferred from a numerical status
string alone.

The old higher-prefix witnesses can be activated only at $t_1\le25/128$.
Both new actions are in $25/128<t_1<19/96$. The second-prefix action
requires $t_2>25/32-3t_1\ge3/16$, whereas the third-prefix action has
$t_2\le247/1728<3/16$. Consequently these are disjoint first-two-prime
cylinders; each new local residual identity has its stated inherited
parent. Neither action changes the other's membership or cap.
The third action keeps its already admitted odd node, and replaces
only its lower residual. The second changes its entire upper residual.
Both preserve the required odd first exits and all selected and full
support gates. Their savings may therefore be added once.

Use the common finite approximation, small-prime separation and signed
compilation from the existing joint audits. The second action retains
its small-third-prime tail argument; the third has $t_3\ge5/96$ and
introduces no such new tail. Trim boundaries, take the old buffers
small enough for the retained histories, and choose the finitely many
new levels below their branch endpoints with strict support margins.
All such choices precede $N\to\infty$. Every $K>E-D_2-D_3$ is thereby
available for the same compiled majorant, its density and its signed
distribution. Endpoint attainment by a single finite tree is not asserted.

Choose
\[
 K=\frac{33809}{5000}=6.7618,\qquad
 G=K-(E-D_2-D_3)>1.0390\cdot10^{-5}.
\]
The full upper-density approximation is allotted $G/4$, and the
remaining compilation and small-prime losses another $G/4$. There is
strict slack after both. The positive fixed-cutoff product theorem in
\texttt{lower-k-uniform-product-interface-review.tex} applies because
its proof uses the same odd-exit, supported-atom, distribution and
density properties just verified, not a particular depth of the old
tree. Its fixed collection of outer prime families is unchanged.
Thus this $K$ can be used both for the final bad-semiprime count and
for the three switched costs below.

\paragraph{Lowering three costs in the existing complete functional.}
The current physical parameters are still
$a=417/5000$, $h=19/125$, $d=59/250$.
Write $K_0=67619/10000$.
In the complete eleven-term functional the three costs have the form
$U_i(K)=KM_i$ for $i=7,8,9$, where the positive masses $M_i$ depend
on these physical parameters but not on the finite $H,h$ table.
Their total contribution is
\[
 -\frac{2U_7(K)+U_8(K)+U_9(K)}4.
\]
The 248-generation baseline replaced their old numerical upper bounds
once by the midpoint certificate. Subsequent pair/four-prefix
refinements, starting-policy changes and restoration increments retained
these three costs at $K_0$. The independent lineage review explicitly
checks this retention.

Let $u_i^-$ be the 112-bit Simpson lower enclosure of $U_i(K_0)$:
\[
 (u_7^-,u_8^-,u_9^-)=2^{-112}
 (7183180737282426411788973366247539,
 21065542644621920270107869035005941,
 16911485600294926597530591281091713).
\]
These are bounds for analytic cost expressions, not lower bounds for
actual prime counts. Since $K_0-K>0$, the functional improves by at
least
\[
 \delta_c=\frac{K_0-K}{4K_0}(2u_7^-+u_8^-+u_9^-)
 >3.72712\cdot10^{-5}.
\]
Keep the other eight terms and the full 248-generation table unchanged.
In particular this does not recalibrate the table, reduce $U_{10},U_{11}$,
or add a second copy of their existing minimum savings. The new upper
tree does not use the Chen coefficient or that table, so this one-way
substitution has no circular premise.

The previously certified endpoint lower bound was
\[
 C_0=\frac{10169845841195726878936621165692227}{2^{112}}.
\]
Retain both old lower-count loss budgets verbatim, each equal to
\[
 L=\frac{828553243746356165863625765747423}
 {64903710731685345356631204115251200000}.
\]
Their finite choices can be kept or refined. Improving these three
positive product bounds uses the already budgeted upper coefficient
$K$, and introduces no larger lower-policy approximation loss. The
finite union of the three product families and the final bad-count
family can use common sufficiently small parameters. Exact arithmetic
gives
\[
 C_0+\delta_c-2L>1.9586528029>
 c:=\frac{39173}{20000}=1.95865.
\]
Hence the existing lower-count argument with the new three costs has
$R(N)\ge(c+o(1))H_N$ under the same analytic inputs.

\paragraph{The imbalance exponent.}
Take $b=748517/1000000$ and $z=(1-b)/(1+b)$. The rational upper bound
\[
 \ell^+=2\sum_{j=0}^{31}\frac{z^{2j+1}}{2j+1}
       +\frac{2z^{65}}{65(1-z^2)}>\log(1/b)
\]
yields
\[
 c-K\ell^+>1.77937\cdot10^{-5}>\frac1{100000}.
\]
The positive bad-count assembly therefore gives at least $H_N/100000$
representations, for every sufficiently large even $N$, with
\[
 N=p+rq,\qquad p,q\text{ prime},\quad r=1\text{ or prime},
 \qquad r\le q^{748517/1000000}.
\]
This is a small improvement of this route's selected safe parameters.
It asserts neither a best known exponent, a new external distribution
theorem, an effective threshold, nor machine formalization. The finer
alternative second-prefix integral, sixth-tail pilots, mixed-first
gain, remaining-$W$ bound and rank-seven diagnostic are not included.
